% Generated draft: requires mathematical and source review.
\documentclass[11pt]{article}
\usepackage[margin=25mm]{geometry}
\usepackage{fontspec}
\setmainfont{Latin Modern Roman}
\usepackage{amsmath,amssymb,mathtools,unicode-math}
\setmathfont{Latin Modern Math}
\usepackage{enumitem,needspace,xurl,hyperref}
\setlength{\parindent}{0pt}
\setlength{\parskip}{5pt}
\setlength{\emergencystretch}{3em}
\newcommand{\N}{\mathbb{N}}
\newcommand{\Z}{\mathbb{Z}}
\newcommand{\Q}{\mathbb{Q}}
\newcommand{\R}{\mathbb{R}}
\newcommand{\C}{\mathbb{C}}
\newcommand{\F}{\mathbb{F}}
\newcommand{\A}{\mathbb{A}}
\newcommand{\PP}{\mathbb{P}}
\newcommand{\E}{\mathbb{E}}
\newcommand{\eps}{\varepsilon}
\newcommand{\dd}{\,\mathrm{d}}
\begin{document}
\section*{UnsolvedMath ID: 30004308}
\textbf{Source problem code:} catalogue entry
\subsection{Definitions and mathematical statement}
Write a set-theoretic Yang--Baxter map as \(r\colon X\times X\to X\times X\), with \(r(x,y)=(\lambda_x(y),\rho_y(x))\). The Yang--Baxter equation is
\[
(r\times\operatorname{id}_X)(\operatorname{id}_X\times r)(r\times\operatorname{id}_X)
=
(\operatorname{id}_X\times r)(r\times\operatorname{id}_X)(\operatorname{id}_X\times r)
\quad\text{on }X^3.
\]
Left nondegeneracy means that \(\lambda_x\colon X\to X\) is bijective for every \(x\in X\). The supplied record does not explicitly specify whether bijectivity of \(r\) is included in its convention for a solution. No finiteness or involutivity hypothesis is stated.

For terminology, a left rack is a set \(R\) with an operation \(\triangleright\) such that each map \(b\mapsto a\triangleright b\) is bijective and
\[
a\triangleright(b\triangleright c)
=(a\triangleright b)\triangleright(a\triangleright c)
\qquad(a,b,c\in R).
\]
In this explicitly chosen notation, its structure group is
\[
G(R)=\left\langle g_a\ (a\in R)\ \middle|\
g_a g_b=g_{a\triangleright b}g_a\quad(a,b\in R)\right\rangle.
\]
This fixes a notation for rack structure groups; the supplied record gives neither a left/right rack convention nor the construction of the derived rack from \(r\). Denote the source's intended derived rack provisionally by \(R_{\mathrm{der}}(X,r)\). The phrase ``abelian rack'' is also undefined in the supplied record and is retained as unresolved terminology.

\textbf{Requested characterization problems.} For left nondegenerate Yang--Baxter solutions in the source's intended class:
\begin{enumerate}
\item Find an intrinsic property of \((X,r)\) equivalent to \(G(R_{\mathrm{der}}(X,r))\) being abelian, that is, to \(uv=vu\) for every \(u,v\) in that group.
\item Find an intrinsic property of \((X,r)\) equivalent to \(R_{\mathrm{der}}(X,r)\) being an abelian rack, after the intended meaning of that term is specified.
\item Determine whether splitting of the source's associated central extension, together with an additional condition, characterizes rack structure groups. Abstractly, a central extension is an exact sequence
\[
1\longrightarrow A\xrightarrow{\iota}E\xrightarrow{\pi}Q\longrightarrow1,
\qquad \iota(A)\subseteq Z(E),
\]
and splitting means that there is a group homomorphism \(s\colon Q\to E\) satisfying \(\pi\circ s=\operatorname{id}_Q\). The record does not identify \(A,E,Q\), their relation to the rack structure group, or the proposed additional condition; consequently, it does not determine a precise equivalence for this question.
\item Give group-theoretic necessary and sufficient conditions for an abstract group \(H\) to be isomorphic to \(G(R)\) for a rack \(R\) in the intended class. The record does not specify whether this class includes all racks or has additional restrictions.
\end{enumerate}
These are characterization requests, not supplied conjectured equivalences. The missing definitions and hypotheses prevent a fully determined formulation of all four questions.
\subsection{Short English statement}
The problem asks how properties of a Yang--Baxter solution determine abelianness of its derived rack and of that rack's structure group, and how to recognize rack structure groups using group theory and central extensions. The derived-rack construction, the meaning of ``abelian rack,'' and the relevant extension and additional condition are missing. This statement-only draft therefore retains those unresolved qualifications for human mathematical review.
\subsection{Sources}
\begin{itemize}
\item Supplied local UnsolvedMath registry record, permanent ID \(30004308\), title ``Abelian Derived Racks and Their Structure Groups,'' collection ``Oberwolfach Reports Open Problems,'' exact source code \texttt{catalogue entry}. The record gives proposed year 2019 and names no proposer.
\item Supplied primary-report citation: \emph{Mini-Workshop: Algebraic Tools for Solving the Yang--Baxter Equation} (2020), \emph{Oberwolfach Reports}, pp.~3207--3242, DOI \texttt{10.4171/owr/2019/51}. \url{https://doi.org/10.4171/owr/2019/51}. The supplied background attributes the problem to workshop open-problem discussions documented in this report; no problem-specific page or speaker attribution is supplied.
\item Supplied literature source: Lebed--Mortier, ``abelian quandles and structure groups.'' \url{https://doi.org/10.1016/j.jpaa.2020.106474}.
\item Supplied literature source: Ced\'o--Jespers--Verwimp, ``derived structure monoids.'' \url{https://arxiv.org/abs/1912.09710}.
\item Supplied literature source labelled ``original OWR report.'' \url{https://doi.org/10.4171/owr/2019/51}.
\item Supplied literature assessment dated \texttt{2026-08-22}: the record labels the problem ``partially solved'' and reports substantial classifications concerning finite rack and quandle structure groups, permutation and derived racks, and abelian-quandle cases, while reporting no general correspondence or group-theoretic characterization answering all four questions. This is supplied dated curation, not an independently verified literature assessment. The background explicitly requests independent verification of time-sensitive status and cited claims.
\end{itemize}
\end{document}
